\section{The minimum-information threshold for positional decoding}
Verdict reframe 2 asks: not ``how accurately can we decode'' but ``what is the minimum biological information required to recover developmental states''. We degrade the Bicoid gradient input to the positional decoder along three axes under the same disjoint-fly-line 5-fold protocol as the head-to-head benchmark, and ask where held-out error crosses the published Liu 2013 true-dose threshold model (4.94\% EL).
\begin{table}[!htbp]\centering\small
\caption{Information-degradation sweep (held-out RMSE, \% egg length). Restricted-input scores are compared with the 4.94\% EL true-dose threshold; this is not the best comparator in the archived benchmark.}\label{tab:infothresh}
\begin{tabular}{lll}\hline\hline
Degradation & level & RMSE (\%EL) \\ \hline
full gradient (64 points) & - & 2.39 \\
anterior window only & $x \le 0.1$ EL & 2.99 \\
sparse sampling & 2 points & 2.53 \\
sparse sampling & \textbf{1 point} & \textbf{4.29} \\
additive noise & $\sigma = 3\times$ signal std & 4.48 \\
additive noise & $\sigma = 4\times$ signal std & 5.03 (crosses) \\
\hline\hline
benchmark: Liu 2013 true-dose model & & 4.94 \\
\hline\hline\end{tabular}
\end{table}

The single-point decoder reaches 4.29\% EL and the anterior-window decoder 2.99\% EL under the tested line-disjoint protocol. Both beat the 4.94\% EL true-dose threshold comparator, but neither beats the measured-amplitude threshold comparator at 2.8566\% EL in \texttt{results/cf\_decoder.json}. The noise crossing is a property of the selected synthetic corruption and estimator, not a minimum biological information threshold. We did not measure information in bits, the embryo's readout circuit, sensing noise or achievable optimal accuracy. Restricted-input prediction therefore does not establish that living morphogen readout needs no gradient shape. The full empirical sweep remains in \texttt{results/info\_threshold.json}.
